\section{Quản lý Dịch vụ (Service Management)}

\subsection{Cài đặt và Quản lý Trạng thái Nginx}

\textbf{Cách thực hiện:}
Cài đặt \texttt{nginx}. Sau đó, sử dụng lệnh \texttt{systemctl} để lần lượt kiểm tra trạng thái, khởi động và dừng dịch vụ web server này.

\begin{lstlisting}[language=bash]
# Cai dat nginx
sudo apt install nginx -y

# Xem trang thai, khoi dong, dung dich vu
sudo systemctl status nginx
sudo systemctl start nginx
sudo systemctl stop nginx
\end{lstlisting}

\textbf{Kết quả:}
Dịch vụ được cài đặt thành công, chuyển đổi trạng thái giữa \texttt{active (running)} và \texttt{inactive (dead)} ổn định qua \texttt{systemctl}.

\begin{figure}[H]
    \centering
    \begin{subfigure}[b]{0.45\textwidth}
        \centering
        \includegraphics[width=\textwidth]{05/install-nginx.png}
        \caption{Cài đặt nginx}
    \end{subfigure}
    \hfill
    \begin{subfigure}[b]{0.45\textwidth}
        \centering
        \includegraphics[width=\textwidth]{05/nginx-status-initial.png}
        \caption{Trạng thái ban đầu}
    \end{subfigure}
    
    \vspace{0.5cm}
    \begin{subfigure}[b]{0.45\textwidth}
        \centering
        \includegraphics[width=\textwidth]{05/nginx-started.png}
        \caption{Nginx active}
    \end{subfigure}
    \hfill
    \begin{subfigure}[b]{0.45\textwidth}
        \centering
        \includegraphics[width=\textwidth]{05/nginx-stopped.png}
        \caption{Nginx inactive}
    \end{subfigure}
    \caption{Quản lý trạng thái service}
\end{figure}

\subsection{Cấu hình Khởi động cùng hệ thống}

\textbf{Cách thực hiện:}
Tắt và bật tính năng tự động khởi động cùng hệ thống của Nginx. Kiểm tra lại bằng cờ \texttt{is-enabled}.

\begin{lstlisting}[language=bash]
sudo systemctl disable nginx
systemctl is-enabled nginx

sudo systemctl enable nginx
systemctl is-enabled nginx
\end{lstlisting}

\textbf{Kết quả:}
Trạng thái hiển thị \texttt{disabled} rồi chuyển thành \texttt{enabled}, xác nhận Nginx sẽ tự boot khi khởi động máy.

\begin{figure}[H]
    \centering
    \begin{subfigure}[b]{0.48\textwidth}
        \centering
        \includegraphics[width=\textwidth]{05/nginx-disabled.png}
        \caption{Disable Nginx}
    \end{subfigure}
    \hfill
    \begin{subfigure}[b]{0.48\textwidth}
        \centering
        \includegraphics[width=\textwidth]{05/nginx-enabled.png}
        \caption{Enable Nginx}
    \end{subfigure}
    \caption{Cấu hình Autostart}
\end{figure}

\subsection{Kiểm tra và Theo dõi File Log}

\textbf{Cách thực hiện:}
- Dùng \texttt{curl} truy cập \texttt{localhost} để tạo request thành công. Đọc file \texttt{access.log}.
- Đọc file \texttt{error.log} để xem lỗi.
- Chạy \texttt{tail -f} để xem log lỗi theo thời gian thực (realtime) khi gửi request sai.

\begin{lstlisting}[language=bash]
# Tao request va xem access log
curl http://localhost
sudo cat /var/log/nginx/access.log

# Xem error log va tail realtime
sudo cat /var/log/nginx/error.log
sudo tail -f /var/log/nginx/error.log
\end{lstlisting}

\textbf{Kết quả:}
Hệ thống ghi nhận chính xác các truy cập (mã 200) và các lỗi truy cập trang không tồn tại (mã 404). Lệnh \texttt{tail -f} giúp giám sát server hiệu quả.

\begin{figure}[H]
    \centering
    \begin{subfigure}[b]{0.45\textwidth}
        \centering
        \includegraphics[width=\textwidth]{05/curl-localhost.png}
        \caption{Curl trang chủ}
    \end{subfigure}
    \hfill
    \begin{subfigure}[b]{0.45\textwidth}
        \centering
        \includegraphics[width=\textwidth]{05/nginx-access-log.png}
        \caption{Access log ghi nhận truy cập}
    \end{subfigure}
    
    \vspace{0.5cm}
    \begin{subfigure}[b]{0.45\textwidth}
        \centering
        \includegraphics[width=\textwidth]{05/nginx-error-log.png}
        \caption{File Error log}
    \end{subfigure}
    \hfill
    \begin{subfigure}[b]{0.45\textwidth}
        \centering
        \includegraphics[width=\textwidth]{05/nginx-tail-realtime.png}
        \caption{Theo dõi Error log realtime}
    \end{subfigure}
    \caption{Quản lý và giám sát File Log}
\end{figure}
